\chapter{Build, packaging, immagini}
\label{cap:20-build-packaging}

Un control plane che si vuole misurabile deve essere anche \emph{producibile in
varianti confrontabili}. Questo capitolo descrive la catena di build di \nf{},
gli assi lungo cui un binario pu\`o variare, e le decisioni di packaging che il
progetto motiva con misure.

\section{Gli assi di variazione}

Un artefatto di \nf{} \`e determinato da cinque scelte indipendenti:

\begin{center}
\small
\begin{adjustbox}{max width=\textwidth}
\begin{tabular}{@{}lp{8.2cm}@{}}
\toprule
\textbf{Asse} & \textbf{Valori} \\
\midrule
Selezione dei moduli & \code{none}, \code{all}, o una lista
(\cref{cap:08-sistema-moduli}) \\
Tipo di build & JVM (Spring Boot fat jar) oppure immagine nativa GraalVM \\
Livello di ottimizzazione nativa & \code{-Os} (dimensione) o \code{-O3}
(velocit\`a, default) \\
Garbage collector nativo & \code{serial} (default Community) o \code{G1}
(Oracle GraalVM) \\
Limite di memoria del container & assente o fissato \\
\bottomrule
\end{tabular}
\end{adjustbox}
\end{center}

Il \cref{cap:24-footprint} misura l'effetto degli ultimi quattro. Che il
progetto li abbia resi tutti selezionabili da riga di comando \`e la condizione
che rende quella misura possibile.

\section{La build JVM}

Java 25, Spring Boot 4.1, Gradle con toolchain risolta automaticamente.
L'immagine OCI \`e costruita su Distroless con una tecnica che vale la pena
descrivere: il runtime Java non \`e un JRE standard ma uno costruito su misura
con \code{jlink}.

\begin{lstlisting}[language=bash, caption={Estratto dal \code{Dockerfile} del
control plane.}]
RUN jar xf app.jar BOOT-INF/lib BOOT-INF/classes && \
    jdeps --ignore-missing-deps --multi-release 25 --recursive \
      --print-module-deps --class-path 'BOOT-INF/lib/*' \
      BOOT-INF/classes app.jar > modules.txt && \
    jlink --add-modules "$(cat modules.txt),jdk.unsupported,jdk.crypto.ec,jdk.httpserver" \
      --strip-debug --no-header-files --no-man-pages --compress=zip-9 \
      --output /jre
\end{lstlisting}

\code{jdeps} determina i moduli JDK effettivamente raggiunti dall'applicazione,
e \code{jlink} ne assembla un runtime minimale. Tre moduli sono aggiunti a mano,
e il commento spiega perch\'e: \code{jdk.unsupported} (Netty usa
\code{sun.misc.Unsafe}), \code{jdk.crypto.ec} (le suite TLS a curve ellittiche
sono caricate via \code{ServiceLoader}) e \code{jdk.httpserver} (il proxy del
provider container). Sono tre casi in cui l'analisi statica di \code{jdeps} non
vede una dipendenza reale, perch\'e la dipendenza \`e risolta a runtime.

\`E lo stesso problema che la compilazione nativa affronta su scala molto
maggiore, e la sua comparsa gi\`a qui, in forma ridotta, \`e istruttiva.

\subsection{I flag della JVM}

\begin{lstlisting}[language=bash]
"-XX:+UseSerialGC", "-XX:MaxRAMPercentage=70", "-XX:TieredStopAtLevel=1",
"-Xss256k", "-Dspring.jmx.enabled=false", ...
\end{lstlisting}

Ciascuno rinuncia a qualcosa in cambio di impronta. \code{TieredStopAtLevel=1}
ferma la compilazione JIT al primo livello: il codice risultante \`e pi\`u lento
ma il compilatore lavora meno e la CodeHeap resta piccola. \code{-Xss256k}
riduce lo stack per thread da 1\,MB al quarto --- con ventitr\'e thread, come
misurato nel \cref{cap:24-footprint}, sono 17\,MB risparmiati.

Va detto che questa configurazione \`e ottimizzata per l'impronta e non per il
throughput di picco, e che le misure di throughput riportate nel
\cref{cap:24-footprint} sono state raccolte su una configurazione diversa. \`E
un'incoerenza da tenere presente nel leggerle.

\begin{damisurare}
Il costo in throughput di \code{TieredStopAtLevel=1} sul control plane non \`e
stato misurato. \`E l'unico flag della lista il cui effetto potrebbe essere
sostanziale sul percorso caldo.
\todomis{throughput con compilazione JIT completa contro C1 soltanto}
\end{damisurare}

\section{La build nativa}

\subsection{Perch\'e}

Tre motivi, in ordine di importanza per questo progetto.

\emph{Il tempo di avvio.} Da 1{,}585\,s a 0{,}072\,s misurati sul control plane
(\cref{cap:24-footprint}). Per un processo di lunga durata \`e ininfluente; per
la CLI \`e la differenza fra usabile e no; per un esperimento che avvia e ferma
il control plane decine di volte \`e tempo di ciclo.

\emph{L'impronta a riposo.} Da 133\,MiB a 25{,}6\,MiB su una funzione Spring.

\emph{L'eliminazione della fase di riscaldamento.} Un'immagine nativa \`e
compilata in anticipo e non ha una curva di riscaldamento JIT. Per un esperimento
questo elimina una variabile: non c'\`e da decidere quanto carico scartare prima
di iniziare a misurare.

\subsection{Il prezzo: la riflessione}

La compilazione anticipata richiede di conoscere staticamente ci\`o che verr\`a
caricato. Ogni uso della riflessione, ogni proxy dinamico, ogni risorsa caricata
per nome deve essere dichiarato. Il progetto lo fa con classi di \emph{runtime
hints} dedicate:

\begin{center}
\small
\begin{adjustbox}{max width=\textwidth}
\begin{tabular}{@{}ll@{}}
\toprule
\textbf{Classe} & \textbf{Copre} \\
\midrule
\code{CaffeineRuntimeHints} & La cache di idempotenza. \\
\code{VertxRuntimeHints} & Vert.x, trascinato dal client Fabric8. \\
\code{DockerJavaRuntimeHints} & Il client Docker Java. \\
\code{GcMetricsConfiguration} & I binder di metriche del collector. \\
\bottomrule
\end{tabular}
\end{adjustbox}
\end{center}

\`E significativo che tre delle quattro riguardino dipendenze dei \emph{moduli}
e non del nucleo. La compilabilit\`a nativa \`e un vincolo che si propaga alle
dipendenze transitive, ed \`e uno degli argomenti a favore della selezione
modulare a tempo di build (\cref{cap:08-sistema-moduli}): un binario che non
include il provider Kubernetes non ha bisogno degli hint di Vert.x.

Il rischio caratteristico \`e che un hint mancante non fa fallire la
compilazione: produce un binario che fallisce a runtime, sul primo percorso che
usa quella riflessione. \`E la ragione per cui esiste lo smoke test nativo della
CLI (\cref{cap:19-cli-scaffolding}).

\subsection{Il livello di ottimizzazione}

Il progetto usa \code{-O3} per default, e il commento nel file di build riporta
la misura che ha motivato il cambio da \code{-Os}:

\begin{center}
\small
\begin{adjustbox}{max width=\textwidth}
\begin{tabular}{@{}lrr@{}}
\toprule
 & \textbf{\code{-Os}} & \textbf{\code{-O3}} \\
\midrule
throughput & 8\,398 rps & \textbf{9\,251 rps} $(+10{,}2\%)$ \\
p95 & 15{,}64\,ms & 14{,}79\,ms \\
massimo & 148{,}7\,ms & 127{,}2\,ms \\
avvio & 0{,}173\,s & \textbf{0{,}068\,s} \\
residente a riposo & 120{,}6\,MiB & \textbf{39{,}3\,MiB} \\
dimensione immagine & 195\,MB & 397\,MB \\
\bottomrule
\end{tabular}
\end{adjustbox}
\end{center}

Il risultato \`e controintuitivo: il livello che ottimizza per la \emph{velocit\`a}
produce anche un'impronta a riposo tre volte pi\`u piccola, a fronte di
un'immagine doppia. Il \cref{cap:24-footprint} discute il meccanismo ---
l'inlining --- e la condizione sotto cui il guadagno svanisce.

Il commento nel build file chiude con l'avvertenza corretta: \emph{``The gain is
conditional [\ldots] under a 1GB limit the two were identical at $\sim$3\,490
requests per second, because there the bottleneck is collection rather than code
quality.''}

\subsection{Il garbage collector}

\begin{lstlisting}[language=Java]
// -PnativeGc=G1 needs Oracle GraalVM on Linux; Community offers only
// 'serial' and 'epsilon', and rejects anything else at build time.
// Measured motivation: the serial collector spent 32% of wall-clock
// time collecting under sustained load, with pauses reaching 1.9s.
if (project.hasProperty('nativeGc')) {
    binary.buildArgs.add("--gc=${project.property('nativeGc')}")
}
\end{lstlisting}

\`E la decisione di packaging pi\`u consequenziale del progetto, ed \`e trattata
in dettaglio nel \cref{cap:24-footprint}. In sintesi: GraalVM Community offre
solo un collector seriale, che su questo carico spende il 32\% del tempo di
parete a raccogliere; G1, disponibile in Oracle GraalVM, elimina il problema. La
selezione resta opt-in (\code{NATIVE\_GC=G1}) perch\'e Oracle GraalVM ha una
licenza diversa.

\subsection{Il budget di memoria del compilatore}

\begin{lstlisting}[language=Java]
// -PnativeBuildMemory=6g bounds the BUILDER's heap, not the built
// image's. native-image sizes its own heap from the memory it can
// see, which is the whole machine even when most of it is already
// spoken for: on a 12GB VM also running k3s, a control plane and
// Prometheus, it was SIGKILLed by the OOM killer after 9m50s of
// work. A build that dies that way wastes the ten minutes and says
// "exit 137"; one that is told its budget either fits or fails fast.
\end{lstlisting}

L'osservazione riguarda una classe generale di problemi: uno strumento che
dimensiona le proprie risorse dalla memoria \emph{visibile} anzich\'e da quella
\emph{disponibile} si comporta male su una macchina condivisa. Il parametro di
parallelismo esiste per lo stesso motivo, con il compromesso opposto: meno
lavoratori costano tempo di parete e risparmiano memoria di picco.

Entrambi sono decisioni che riguardano una macchina su cui la build compete con
il carico contro cui verr\`a poi misurata --- una situazione specifica di un
progetto sperimentale, e non affrontata dalla documentazione degli strumenti.

\section{Una dipendenza esclusa}

\begin{lstlisting}[language=Java]
configurations.all {
    // HTTP/3/QUIC support that reactor-netty pulls in automatically; nothing in this
    // codebase enables HTTP/3. Drops netty-codec-http3 plus its 5 platform-native quic jars.
    exclude group: 'io.netty', module: 'netty-codec-http3'
}
\end{lstlisting}

Sei archivi rimossi, di cui cinque contenenti codice nativo per piattaforma
specifica. Nel contesto della compilazione nativa e delle immagini Distroless,
una dipendenza inutilizzata che porta librerie native non \`e soltanto peso: \`e
superficie su cui gli hint di runtime potrebbero servire e non ci sono.

Il progetto utilizza il plugin \code{dependency-analysis} per individuare
sistematicamente dipendenze non usate o non dichiarate
(\cref{cap:26-qualita-verifica}).

\section{La matrice degli artefatti}

\begin{center}
\small
\begin{adjustbox}{max width=\textwidth}
\begin{tabular}{@{}lll@{}}
\toprule
\textbf{Artefatto} & \textbf{Build JVM} & \textbf{Build nativa} \\
\midrule
control plane & \code{Dockerfile} + Distroless & \code{native-java-image.sh control-plane} \\
\code{warm-echo} & \code{Dockerfile} + Distroless & \code{native-java-image.sh warm-echo} \\
funzioni Java & per funzione & \code{native-java-image.sh <nome>} \\
CLI & \code{installDist} & \code{nativeCompile} \\
watchdog & --- & Rust statico musl, \code{FROM scratch} \\
\bottomrule
\end{tabular}
\end{adjustbox}
\end{center}

Un unico comando costruisce e verifica l'intera matrice nativa:

\begin{lstlisting}[language=bash]
scripts/native-build.sh
\end{lstlisting}

Le immagini native non usano i buildpack di Spring Boot ma un
\code{Dockerfile} che copia il binario su Distroless. La scelta \`e coerente con
il resto: un buildpack decide da s\'e i flag di compilazione, e questo progetto
ha bisogno che quei flag siano una variabile sperimentale esplicita.
